%!TEX program = lualatex
% Arquivo: relatorio_pratica_3_HCl_combinado.tex

\documentclass[a4paper,12pt]{article}

% =========================================================
% Idioma e tipografia (LuaLaTeX)
% =========================================================
\usepackage[brazilian]{babel}
\usepackage{fontspec}
\setmainfont{Latin Modern Roman}
\usepackage{microtype}

% =========================================================
% Layout
% =========================================================
\usepackage[top=3cm,bottom=2.5cm,left=3cm,right=2.5cm]{geometry}
\usepackage{setspace}
\onehalfspacing
\usepackage{indentfirst}
\setlength{\parindent}{1.25cm}

% =========================================================
% Matemática, unidades e fórmulas químicas
% =========================================================
\usepackage{amsmath}
\usepackage{amssymb}
\usepackage{siunitx}
\sisetup{
  locale = DE,
  output-decimal-marker = {,},
  per-mode = symbol,
  exponent-product = \cdot
}
\usepackage[version=4]{mhchem}

% =========================================================
% Tabelas, listas, figuras e desenhos
% =========================================================
\usepackage{booktabs}
\usepackage{array}
\usepackage{enumitem}
\setlist[itemize]{noitemsep,topsep=3pt}
\setlist[enumerate]{noitemsep,topsep=3pt}

\usepackage{graphicx}
\usepackage{float}
\usepackage{pdflscape}
\graphicspath{{./}{./imagens/}{./figuras/}}

\usepackage{xcolor}
\usepackage{tikz}
\usetikzlibrary{arrows.meta,calc,positioning,shapes.geometric}
\usepackage[most]{tcolorbox}

% =========================================================
% Links
% =========================================================
\usepackage{hyperref}
\usepackage{xurl}
\hypersetup{
  colorlinks=true,
  linkcolor=black,
  urlcolor=blue,
  citecolor=black,
  pdfauthor={Israel Carneiro; José Mauro Xavier Elsas; Miguel Gayani},
  pdftitle={Prática 3 -- Preparo e padronização de solução de HCl}
}

% =========================================================
% Cores e comandos usados nos esquemas da prática
% =========================================================
\definecolor{quadroazul}{RGB}{40,91,140}
\definecolor{quadrovermelho}{RGB}{177,69,60}
\definecolor{quadrocinza}{RGB}{245,245,245}

\newcommand{\subtitulo}[1]{%
  \par\medskip
  {\color{quadroazul}\bfseries\large #1}\par\smallskip
}

% =========================================================
% Dados da capa
% =========================================================
\newcommand{\Instituicao}{Universidade do Estado do Rio de Janeiro -- UERJ}
\newcommand{\Curso}{Química VIII / Oceanografia}
\newcommand{\Professor}{Professor: Ítalo Vieira}
\newcommand{\Aluno}{Alunos: Israel Carneiro\\
José Mauro Xavier Elsas\\
Miguel Gayani}
\newcommand{\Cidade}{Rio de Janeiro}
\newcommand{\Ano}{2026}

\title{\textbf{Prática 3}}
\author{\Aluno}
\date{\Cidade{} -- \Ano}

\begin{document}

% =========================================================
% Capa (página de rosto)
% =========================================================
\begin{titlepage}
\centering
{\large \Instituicao \par}
\vspace{0.5cm}
{\large \Curso \par}
\vspace{1.0cm}

{\LARGE \textbf{Prática 3} \par}
\vspace{0.6cm}
{\Large \textbf{PREPARO E PADRONIZAÇÃO DE SOLUÇÃO\\
DE \ce{HCl} \qty{0.1}{\mole\per\litre}} \par}

\vspace{1cm}

% O logotipo será inserido quando o arquivo uerj.jpg estiver na pasta
% principal, em imagens/ ou em figuras/.
\IfFileExists{uerj.jpg}{%
  \begin{figure}[H]
    \centering
    \includegraphics[width=0.3\linewidth]{uerj.jpg}
  \end{figure}
}{%
  \begin{tcolorbox}[
    width=0.32\textwidth,
    colback=white,
    colframe=black!35,
    boxrule=0.5pt,
    arc=0mm
  ]
  \centering
  \vspace{1.2cm}
  \textit{Logotipo da UERJ}
  \vspace{1.2cm}
  \end{tcolorbox}
}

\vfill

\begin{flushright}
\Aluno\\
\Professor
\end{flushright}

\vspace{1.0cm}
{\Cidade{} -- \Ano \par}
\end{titlepage}

% =========================================================
% Sumário
% =========================================================
% \tableofcontents
% \newpage

\section{Introdução}

Os reagentes mais comumente usados para o preparo de soluções padrão de ácidos são o ácido
clorídrico, o ácido sulfúrico e o ácido perclórico. A solução de ácido clorídrico é frequentemente preferida, pois:

\begin{itemize}
    \item é estável por longos períodos;
    \item pode ser usada na presença da maior parte dos cátions sem sofrer interferência significativa decorrente da formação de sais pouco solúveis.
\end{itemize}

A solução concentrada de \ce{HCl} apresenta, aproximadamente, \qty{37}{\percent} em massa e densidade de \qty{1.19}{\gram\per\milli\litre}. Como essa solução não constitui um padrão primário, a solução de \ce{HCl} \qty{0.1}{\mole\per\litre} deve ser inicialmente preparada por diluição de um volume calculado do ácido concentrado e, posteriormente, padronizada.

\section{Material Utilizado}

\begin{enumerate}[label=\alph*.]
    \item cilindro graduado de \qty{500}{\milli\litre};
    \item balão volumétrico de \qty{500}{\milli\litre};
    \item balão volumétrico de \qty{250}{\milli\litre};
    \item pipeta graduada;
    \item béquer de \qtyrange{200}{300}{\milli\litre};
    \item bastão de vidro;
    \item funil;
    \item dois frascos de vidro ou plástico com capacidade de \qty{500}{\milli\litre}.
\end{enumerate}

\section{Reagentes}

\begin{enumerate}[label=\alph*.]
    \item ácido clorídrico concentrado, \ce{HCl}, grau p.a.;
    \item solução do indicador alaranjado de metila;
    \item tetraborato de sódio deca-hidratado, \ce{Na2B4O7.10H2O} (bórax), grau p.a.
\end{enumerate}

\section{Procedimento Experimental}

\subsection{Preparo da solução de \ce{HCl} \qty{0.1}{\mole\per\litre}}

Para preparar \qty{500}{\milli\litre} de solução de \ce{HCl}
\qty{0.1}{\mole\per\litre}, determinou-se inicialmente a quantidade de matéria de ácido necessária.

\begin{minipage}[t]{0.48\textwidth}
\subtitulo{Número de mols}

\begin{align*}
  n &= C\,V \\
    &= \left(\qty{0.1}{\mole\per\litre}\right)
       \left(\qty{0.5}{\litre}\right) \\
    &= \boxed{\qty{0.05}{\mole}}
\end{align*}

\subtitulo{Massa de \ce{HCl} puro}

\begin{align*}
  m_{\ce{HCl}} &= MM\,n \\
    &= \left(\qty{36.5}{\gram\per\mole}\right)
       \left(\qty{0.05}{\mole}\right) \\
    &= \boxed{\qty{1.825}{\gram}}
\end{align*}

\begin{tcolorbox}[
  title=Etiqueta da solução,
  colback=quadrocinza,
  colframe=black!65,
  boxrule=0.6pt,
  arc=1mm
]
\begin{tabular}{@{}ll@{}}
Data: & \rule{3.0cm}{0.4pt}\\[1.3ex]
Nome da solução: & \rule{2.2cm}{0.4pt}\\[1.3ex]
Nome do grupo: & \rule{2.5cm}{0.4pt}\\[1.3ex]
Turma: & 2\\[1.3ex]
Concentração: & \rule{2.8cm}{0.4pt}
\end{tabular}
\end{tcolorbox}
\end{minipage}
\hfill
\begin{minipage}[t]{0.48\textwidth}
\begin{tcolorbox}[
  title=Dados da solução concentrada de \ce{HCl},
  colback=quadrocinza,
  colframe=black!65,
  boxrule=0.6pt,
  arc=1mm
]
\[
  T = \qty{37}{\percent}\;(m/m)
  \qquad
  d = \qty{1.19}{\gram\per\milli\litre}
\]
\end{tcolorbox}

\subtitulo{Massa da solução concentrada}

Como a solução comercial contém \qty{37}{\percent} de \ce{HCl} em massa:

\begin{align*}
  m_{\mathrm{sol.\,conc.}}
    &= m_{\ce{HCl}}\frac{100}{T(\%)} \\
    &= \qty{1.825}{\gram}\frac{100}{37} \\
    &= \boxed{\qty{4.93}{\gram}}
\end{align*}

\subtitulo{Volume da solução concentrada}

\begin{align*}
  V_{\mathrm{sol.\,conc.}}
    &= \frac{m_{\mathrm{sol.\,conc.}}}{d} \\
    &= \frac{\qty{4.932}{\gram}}
             {\qty{1.19}{\gram\per\milli\litre}} \\
    &\approx \boxed{\qty{4.1}{\milli\litre}}
\end{align*}
\end{minipage}

\vspace{0.5cm}

\begin{tcolorbox}[
  colback=quadrovermelho!6,
  colframe=quadrovermelho,
  boxrule=1pt,
  arc=2mm
]
\centering
{\color{quadrovermelho}\bfseries\Large Adicionar sempre o ácido sobre a água.}
\end{tcolorbox}

\begin{figure}[H]
\centering
\begin{tikzpicture}[>=Latex,thick,font=\small]
  % Pipeta
  \draw[quadrovermelho,fill=quadrovermelho!8]
    (-3.6,3.4) -- (-3.25,3.4) -- (-3.25,1.65)
    -- (-3.38,1.25) -- (-3.51,1.65) -- cycle;
  \draw[quadrovermelho] (-3.62,3.55) -- (-3.23,3.55);
  \node[quadrovermelho,above] at (-3.42,3.55) {pipeta};
  \node[quadrovermelho,align=center,left=3mm] at (-3.65,2.35)
    {\qty{4.1}{\milli\litre}\\de \ce{HCl} concentrado};

  % Seta para o cilindro
  \draw[->,quadrovermelho] (-3.05,1.55) .. controls (-2.4,1.1) and (-1.8,1.0) .. (-1.1,1.25);

  % Cilindro graduado
  \draw[quadrovermelho,fill=white] (-0.8,0) rectangle (0.2,3.4);
  \draw[quadrovermelho] (-0.95,3.4) -- (0.35,3.4);
  \draw[quadrovermelho] (-1.0,0) -- (0.4,0);
  \fill[quadroazul!16] (-0.77,0.05) rectangle (0.17,1.75);
  \draw[quadroazul] (-0.77,1.75) -- (0.17,1.75);
  \foreach \y in {0.4,0.8,1.2,1.6,2.0,2.4,2.8,3.2}
    \draw[quadrovermelho] (-0.8,\y) -- (-0.55,\y);
  \node[quadrovermelho,below=2mm] at (-0.3,0) {cilindro de \qty{500}{\milli\litre}};
  \draw[->,quadrovermelho] (2.8,2.05) -- (0.3,1.78);
  \node[quadrovermelho,align=left,right] at (2.85,2.05)
    {\qty{250}{\milli\litre} de\\água destilada};

  % Balão volumétrico final
  \draw[->,black!65] (0.65,0.9) -- (2.15,0.9);
  \draw[black!70]
    (2.65,0.25) .. controls (2.05,0.65) and (2.15,1.45) .. (2.55,1.75)
    -- (2.55,3.05) -- (2.95,3.05) -- (2.95,1.75)
    .. controls (3.35,1.45) and (3.45,0.65) .. (2.85,0.25)
    .. controls (2.78,0.20) and (2.72,0.20) .. (2.65,0.25);
  \draw[black!70] (2.48,3.05) -- (3.02,3.05);
  \draw[quadroazul] (2.53,1.55) -- (2.97,1.55);
  \node[below=2mm,align=center] at (2.75,0.22)
    {balão volumétrico\\de \qty{500}{\milli\litre}};
  \node[align=center,right=4mm] at (3.4,1.2)
    {transferir e completar\\até o traço de aferição};
\end{tikzpicture}
\caption{Esquema do preparo da solução de \ce{HCl} \qty{0.1}{\mole\per\litre}.}
\label{fig:preparo_hcl}
\end{figure}

O procedimento consistiu em adicionar inicialmente cerca de \qty{250}{\milli\litre} de água destilada ao recipiente de preparo. Em seguida, mediram-se cuidadosamente \qty{4.1}{\milli\litre} de \ce{HCl} concentrado, que foram adicionados lentamente sobre a água. A solução foi transferida quantitativamente para um balão volumétrico de \qty{500}{\milli\litre}, completando-se o volume com água destilada até o traço de aferição. Por fim, a solução foi homogeneizada e identificada.

\begin{tcolorbox}[
  title=Observação técnica,
  colback=quadrocinza,
  colframe=black!65,
  boxrule=0.6pt,
  arc=1mm
]
Na anotação do quadro, o recipiente intermediário foi identificado como ``cilindro de \qty{500}{\milli\litre}''. Para obter com maior precisão o volume final da solução, o ajuste a \qty{500}{\milli\litre} deve ser realizado em balão volumétrico.
\end{tcolorbox}

\subsection{Preparo da solução de bórax \qty{0.025}{\mole\per\litre}}

Para preparar \qty{250}{\milli\litre} de solução de bórax, considerando a massa molar do tetraborato de sódio deca-hidratado igual a \qty{381.37}{\gram\per\mole}, calculou-se a massa necessária pela relação:

\begin{align*}
  M &= \frac{m}{MM\,V}
  &\Longrightarrow\qquad
  m &= M\,MM\,V \\
  m &= \left(\qty{0.025}{\mole\per\litre}\right)
       \left(\qty{381.37}{\gram\per\mole}\right)
       \left(\qty{0.250}{\litre}\right)
  &= \boxed{\qty{2.38}{\gram}}
\end{align*}

\begin{figure}[H]
\centering
\begin{tikzpicture}[>=Latex,thick,font=\small,node distance=2cm]
  % Béquer 1
  \begin{scope}[shift={(-5,0)}]
    \draw[quadrovermelho] (-0.75,1.8) -- (-0.55,0.15) -- (0.55,0.15) -- (0.75,1.8);
    \draw[quadrovermelho] (-0.85,1.8) -- (0.85,1.8);
    \fill[quadrovermelho!15] (-0.53,0.18) rectangle (0.53,0.45);
    \node[quadrovermelho,below=2mm] at (0,0.15) {béquer};
    \node[quadrovermelho,align=center,above=4mm] at (0,1.8)
      {pesar \qty{2.38}{\gram}\\de bórax};
  \end{scope}

  % Seta 1
  \draw[->,quadrovermelho] (-3.7,1.0) -- (-1.7,1.0);

  % Béquer 2 com água e bastão
  \begin{scope}[shift={(0,0)}]
    \draw[quadrovermelho] (-0.75,1.8) -- (-0.55,0.15) -- (0.55,0.15) -- (0.75,1.8);
    \draw[quadrovermelho] (-0.85,1.8) -- (0.85,1.8);
    \fill[quadroazul!16] (-0.61,0.18) -- (-0.50,1.15) -- (0.50,1.15) -- (0.61,0.18) -- cycle;
    \draw[quadroazul] (-0.50,1.15) -- (0.50,1.15);
    \draw[black!65] (-0.15,0.4) -- (0.28,2.35);
    \node[quadrovermelho,align=center,above right=1mm] at (0.8,1.4)
      {\qty{150}{\milli\litre} de\\água destilada};
    \node[below=2mm,align=center] at (0,0.15)
      {dissolver e homogeneizar};
  \end{scope}

  % Seta 2
  \draw[->,quadrovermelho] (1.7,1.0) -- (3.4,1.0);

  % Balão volumétrico
  \begin{scope}[shift={(5,0)}]
    \draw[quadrovermelho]
      (-0.22,2.55) -- (-0.22,1.65)
      .. controls (-0.95,1.20) and (-0.90,0.35) .. (-0.30,0.08)
      .. controls (-0.10,0.00) and (0.10,0.00) .. (0.30,0.08)
      .. controls (0.90,0.35) and (0.95,1.20) .. (0.22,1.65)
      -- (0.22,2.55);
    \draw[quadrovermelho] (-0.33,2.55) -- (0.33,2.55);
    \draw[quadroazul] (-0.22,1.55) -- (0.22,1.55);
    \node[quadrovermelho,below=2mm,align=center] at (0,0.08)
      {balão de \qty{250}{\milli\litre}};
    \node[align=center,right=4mm] at (0.65,1.15)
      {transferir quantitativamente\\e completar até o traço};
  \end{scope}
\end{tikzpicture}
\caption{Esquema do preparo da solução de bórax \qty{0.025}{\mole\per\litre}.}
\label{fig:preparo_borax}
\end{figure}

A massa de \qty{2.38}{\gram} de bórax foi pesada em um béquer. Adicionaram-se aproximadamente \qty{150}{\milli\litre} de água destilada, e a mistura foi agitada com bastão de vidro até a dissolução do sólido. A solução foi transferida quantitativamente para um balão volumétrico de \qty{250}{\milli\litre}. O béquer e o funil foram lavados com pequenas porções de água destilada, que também foram transferidas para o balão. Em seguida, completou-se o volume até o traço de aferição e homogeneizou-se a solução.

% A etapa de padronização pode ser inserida a seguir quando os dados e o
% procedimento experimental estiverem completos.
%
% \subsection{Padronização da solução de \ce{HCl}}

% =========================================================
% Referências
% =========================================================
\newpage
\phantomsection
\addcontentsline{toc}{section}{Referências}
\begin{thebibliography}{99}

\bibitem{uerj_quimica8_2017}
LAGO, Dalva Cristina Baptista do;
GOIS, Jefferson Santos de;
LANGONE, Marta Antunes Pereira;
SENNA, Lilian Ferreira de.
\textbf{Química VIII -- Oceanografia: aulas práticas}.
Rio de Janeiro: Universidade do Estado do Rio de Janeiro, Instituto de Química, Departamento de Química Analítica, 2017.

\end{thebibliography}

\end{document}
